\section{Step 32 Compress the first padding block and restart its count}\label{sha-s32}

Call compression on the now-complete buffer, then clear the working used count. Keep the saved original bit length in \texttt{s3} across the call. The same physical buffer will hold the final block next; it is its logical occupancy that has become zero. Its old bytes still need to be overwritten where the final block requires zeros.

\subsection{Write and inspect}
Replace only the temporary instructions between \texttt{\# BEGIN fin\_first} and \texttt{\# END fin\_first}. Keep the supplied labels outside that region.

\Needspace{10\baselineskip}
\MARSCodeLabel{MIPS assembly}
\begin{minted}[samepage=true]{mars-idea}
    move $a0, $s0
    daddiu $a1, $s0, 48
    daddiu $a2, $s0, 112
    jal sha256_compress
    nop
    move $s2, $zero
\end{minted}

Use recipe \texttt{s32-1} with the cumulative program through this edit. The exact recipe names the entry, prepared input and checkpoint. Observe the checkpoint before continuing to its temporary or normal return.

\subsection{Expected observation}
The state includes the first padding block and \texttt{s2} is zero before the final-block fill loop.

\Needspace{9\baselineskip}
\subsection{Check your understanding}
May the final block reuse the same 64-byte storage?

\begin{itemize}
\item \textbf{A} No, each block needs a separate allocation.
\item \textbf{B} Yes, after the previous block has been consumed.
\item \textbf{C} Only for ASCII input.
\end{itemize}
